\documentclass[11pt,a4paper]{amsart}
\usepackage[T1]{fontenc}
\usepackage{amssymb}
\usepackage[hmargin=1in,vmargin=.9in]{geometry}
\usepackage[hidelinks]{hyperref}
\hypersetup{pdftitle={The complexity of persuasion with a privately informed binary-action receiver}}
\DeclareFontShape{OT1}{cmr}{bx}{n}{<-> cmbx10}{}
\DeclareMathSizes{8}{8}{8}{8}
\DeclareMathSizes{9}{9}{8}{8}
\DeclareMathSizes{10.95}{10.95}{8}{8}
\newtheorem{theorem}{Theorem}
\title{The complexity of persuasion with a privately informed binary-action receiver}
\date{}
\begin{document}
\begin{abstract}
Computing a payoff-maximizing offline signaling scheme for a privately
informed binary-action receiver is NP-hard when the private channel is
known. With rational data and sender-favorable ties, the payoff-threshold
decision problem is NP-complete, even when the sender's utility is one
for action~1 and zero for action~0 in every state. A reduction from
3-SAT uses a rare anchor state to reward satisfied tests and paired
states whose imbalance incurs a larger payoff loss through guard
signals. It creates an inverse-polynomial additive payoff gap and rules
out a deterministic additive FPTAS unless \(\mathrm P=\mathrm{NP}\).
\end{abstract}
\maketitle

\section{Offline persuasion problem}

Let \(\Theta\) be a finite state set with a full-support prior \(\mu\).
The sender commits to a channel \(\phi(s\mid\theta)\) for messages \(s\).
The receiver privately sees a signal \(r\) from a known finite
channel \(\psi(r\mid\theta)\), conditionally independent of \(s\)
given \(\theta\), and chooses action 0 or 1. The two payoff functions
are \(u_R(j,\theta),u_S(j,\theta)\in[0,1]\). The receiver maximizes
posterior receiver utility and chooses a sender-preferred action when
receiver utility is tied. All input numbers and the payoff threshold
are binary-encoded rationals. The private alphabet is part of the input.

We use the offline formulation of Vurankaya and Topcu
\cite[eqs.~(2)--(3)]{vt}. The question after their Lemma~6 asks for
polynomial-time computation with two receiver actions. Their online
estimation algorithm uses additional regularity and margin conditions.
Here the true private channel is supplied.

\begin{theorem}
Deciding whether an offline instance reaches a supplied sender payoff
threshold is NP-complete. Hardness holds when \(u_S(j,\theta)=j\).
Moreover, a deterministic fully polynomial additive approximation scheme
for the optimal value in this restriction would imply
\(\mathrm P=\mathrm{NP}\).
\end{theorem}

\section{Rational certificates}

Write \(q_s(\theta)=\mu(\theta)\phi(s\mid\theta)\), so \(q_s\geq0\)
and \(\sum_s q_s=\mu\).
Let \(\Delta_R(\theta)=u_R(1,\theta)-u_R(0,\theta)\).
The unnormalized receiver score on \(s,r\) is
\[
 D_r(q_s)=\sum_\theta q_s(\theta)\psi(r\mid\theta)\Delta_R(\theta).
\]
Action 1 is a best response for \(D_r\geq0\), and action 0 for
\(D_r\leq0\). At zero score, the sender-favorable tie rule compares
the two unnormalized sender utilities. At zero probability both
utilities contribute zero.

A proposed binary action profile \(b=(b_r)_r\) for a message imposes
\((2b_r-1)D_r(q)\geq0\) for every \(r\). Its assigned payoff
\[
 L_b(q)=\sum_{\theta,r}q(\theta)\psi(r\mid\theta)u_S(b_r,\theta)
\]
is linear and no greater than the actual payoff under the
tie rule whenever \(q\) satisfies those constraints.
Normalize any nonzero \(q\) to lie in the simplex. For each profile
its weak best-response inequalities cut out a bounded rational
polytope. Its vertices have polynomial bit length by the
determinant bounds for rational linear systems.

Take the finite linear program whose columns are these vertices,
with one column for each compatible profile. Its variable for a
column is a probability, its only equality constraints say the
weighted state vector is \(\mu\), and its objective uses the
assigned column payoffs. Any signaling policy can use its actual
profiles and decompose each posterior into such vertices, preserving
Bayes plausibility and the assigned payoff. More generally, the finite
convex hull of normalized state-and-payoff vectors is closed,
so the same assertion holds for a distribution of posteriors.
Conversely, each feasible linear-program solution is a feasible
signal distribution and its actual payoff is at least the objective.
The optima therefore agree and are attained.

There is an optimal basic solution with at most \(|\Theta|\) positive
weights, again rational with polynomial bit length by determinant
bounds. The certificate lists its corresponding mass vectors and
profiles. A verifier checks nonnegativity, prior sums, weak
incentives and the assigned payoff. These operations take polynomial
bit time. This proves membership in NP for general input payoffs.
For fixed private alphabet size \(k\), one can instead group messages
by profile and keep one unnormalized mass vector for each of the
\(2^k\) profiles. This gives a polynomial-sized rational linear
program when \(k\) is fixed.

\section{Tests for a Boolean formula}

Let a 3-CNF have \(n\geq1\) Boolean variables and \(m\geq1\) clauses.
Use nonnegative abstract variables \(t,x_1,\ldots,x_n\) and an anchor
mass \(z\). For each variable include the mandatory form \(t-x_i\)
and the two alternatives \(t/8-x_i\) and \(x_i-7t/8\).
A clause adds the sum of its three literal values minus \(t/2\), with
literal values \(x_i\) for positive occurrences and \(t-x_i\) for
negative occurrences. There are \(M=3n+m\) tests, and set \(K=2n+m\).
Repeated literals are allowed. Every combined form coefficient is
at most 3 in absolute value.

For \(z>0\), at most one alternative per variable can be at least
\(z/2\), because their sum equals \(-3t/4\). If \(K\) tests reach
\(z/2\), every mandatory form, every clause, and one alternative per
variable reaches it. Assign false for the low alternative and true
for the high alternative. The low values satisfy \(x_i\leq t/8\);
the high values satisfy \(7t/8\leq x_i\leq t\).
Under this assignment a false clause has all three literal values at
most \(t/8\). Its form is at most \(-t/8\), a contradiction.
An unsatisfiable formula thus permits at most \(K-1\) passing tests.

Conversely, fix a satisfying assignment and \(z>0\). Put \(t=32z\)
and set \(x_i=t/16\) for false and \(x_i=15t/16\) for true.
All mandatory forms and chosen alternatives have value at least \(t/16\).
Each satisfied clause has one true literal of value \(15t/16\) and
the other literal values are nonnegative. These \(K\) forms all exceed \(z\).

\section{From tests to a private channel}

For each of the \(n+1\) abstract variables use one plus and one minus
state. Add one anchor state with mass
\(\mu_z=1/[256(n+1)]\). Each paired state has prior
\(p=(1-\mu_z)/[2(n+1)]\).
The difference \(\Delta_R\) is \(+1\) on plus states, \(-1\) on
minus states, and \(-\delta\) on the anchor, where
\[
 a=\frac{1}{64M},\qquad b=\frac12,\qquad
 \delta=\frac{\mu_z}{12M}.
\]
Taking positive and negative parts gives receiver payoffs in
\([0,1]\). The sender gets one for action 1 and zero for action 0.

A private guard signal for a pair has likelihood \(b\) on both
states of the pair and zero elsewhere. For each test form, a test
signal has likelihood \(a\) on the anchor, and assigns \(a\delta\)
times the absolute value of each nonzero coefficient to the plus
copy for a positive coefficient or the minus copy for a negative
coefficient. Unassigned likelihoods are zero.
On a paired state the used likelihood is at most
\(b+3Ma\delta<1\). On the anchor it is \(Ma=1/64<1\).
A distinct singleton signal for each state fills its remaining
likelihood, so the channel is stochastic. The singleton payoff
\(P_{\mathrm{pad}}\) is a computable constant. A singleton
reveals its state, and only plus states induce action 1.

For a message with anchor mass \(z\), guard \(i\) induces action 1
exactly when \(q_i^+\geq q_i^-\). The tie rule selects action 1 at
equal masses. A test induces action 1 exactly when
\(L_{\mathrm{split}}\geq z\), where \(L_{\mathrm{split}}\) evaluates
each coefficient using its assigned sign copy. Its score has the
positive multiplier \(a\delta\). The maximum total guard payoff
is \(b(1-\mu_z)\). Test payoffs coming from all paired states are
at most \(3Ma\delta=a\mu_z/4\), by summing their likelihood bounds
at the prior. Let \(C_0=P_{\mathrm{pad}}+b(1-\mu_z)\).

\section{Completeness, soundness, and gap}

Suppose the formula is satisfiable. Use a witness message with
anchor mass \(\mu_z\) and equal paired masses given by the satisfying
construction. Every paired coordinate is at most \(32\mu_z<p\).
A second message receives all residual mass, so its paired
coordinates are equal and its anchor mass is zero. All guards
succeed and at least \(K\) tests reward the anchor. Hence
\[
 V_{\mathrm{yes}}\geq C_0+a\mu_z K.
\]

Now suppose the formula is unsatisfiable. For an arbitrary finite
scheme define
\[
 F=\sum_{s,i:q_{s,i}^+<q_{s,i}^-}q_{s,i}^- .
\]
Failed guards lose at least \(bF\) in total. Put
\(\Delta_s=\sum_i|q_{s,i}^+-q_{s,i}^-|\) and \(I=\sum_s\Delta_s\).
Equal paired priors imply that the total positive discrepancy equals
the absolute total negative discrepancy. The latter is at most \(F\),
so \(I\leq2F\).

At a message evaluate each abstract variable at the mean of its
two masses. Each sign copy differs from the mean by half the
absolute imbalance, and the coefficients have magnitude at most
3, so \(|L_{\mathrm{split}}-L|\leq3\Delta_s\).
For \(z_s>0\) and \(3\Delta_s\leq z_s/2\), a successful test must
have \(L\geq z_s/2\). At most \(K-1\) tests can succeed in such
a message. The remaining anchor-bearing messages have total
anchor mass at most \(6I\leq12F\). Even if all \(M\) of their tests
succeed, they yield no more than \(12MaF\) anchor payoff.
Adding the padding, guards and paired-state test bound gives
\[
 \begin{split}
 V_{\mathrm{no}}
 &\leq C_0-bF+a\mu_z(K-1)+12MaF+a\mu_z/4\\
 &\leq C_0+a\mu_z(K-3/4),
 \end{split}
\]
since \(12Ma=3/16<b\). The argument covers the optimal finite
scheme proved to exist above.

Set the decision threshold at \(C_0+a\mu_z(K-1/2)\).
The yes and no cases are separated from it by at least \(a\mu_z/4\).
There are \(2(n+1)+1\) states and \(M+3(n+1)+1\) private signals.
The likelihood table, payoffs, and threshold are rational with
polynomial encoding lengths. This is a polynomial-time reduction
from 3-SAT. An exact optimal scheme in explicit rational form can be
evaluated by receiver signs in polynomial bit time.
An additive FPTAS, using error less than \(a\mu_z/4\), would
decide 3-SAT in polynomial time because \(1/(a\mu_z)\) is polynomial
in the formula size. This proves the theorem.

\section{Prior work and validation}

Hossain et al. \cite[Thm.~3 and App.~A.3]{hossain} prove earlier
best-response hardness with \(2k+1\) actions in their reduction.
Dughmi and Xu \cite[Thm.~6.2]{dx} prove public-persuasion hardness for
multiple binary receivers and average adoption. Castiglioni et al.
\cite[Thm.~1]{cast} prove offline hardness with type-dependent
receiver utilities. In the known-channel model, however, every
nonzero coefficient \(\psi(r\mid\theta)\Delta_R(\theta)\) at
a fixed state has the same sign. The present construction enforces
this restriction through sign copies and controls their imbalance
with guards. Sender-favorable ties are used at balanced guards.
The separate online learning question is outside the result.

Internally reviewed, not peer reviewed. Exact rational replay checks
the constructed channels, witness messages, and payoff bounds in
small worked examples. The partial Lean 4.30.0 certificate checks rational test inequalities,
the small-imbalance transfer and the final payoff comparison. The message
optimization and complexity proof are analytic.

Under human direction, ChatGPT Astra 6.2 (OpenAI) generated the
discovery proof and its research packet. OpenAI Codex independently checked the mathematics.
GPT-6.2 Sol (OpenAI) checked the literature and prepared the
manuscript, exact replay, partial Lean certificate and illustrations.
DannyExperiments maintains the repository and is the preferred
citation identity. No AI system is proposed as an author.

\enlargethispage{3\baselineskip}
\begin{thebibliography}{9}
\raggedright
\bibitem{vt}
I. A. Vurankaya and U. Topcu, \emph{Learning to Persuade
Privately Informed Receivers}, arXiv:2607.28342v1, 2026.
\href{https://arxiv.org/abs/2607.28342}{arXiv:2607.28342}.
\bibitem{hossain}
S. Hossain, T. Wang, T. Lin, Y. Chen, D. C. Parkes and H. Xu,
\emph{Multi-Sender Persuasion: A Computational Perspective},
Proc. ICML, PMLR 235 (2024), 18944--18971.
\href{https://proceedings.mlr.press/v235/hossain24c.html}{PMLR publication}.
\bibitem{dx}
S. Dughmi and H. Xu, \emph{Algorithmic Persuasion with No Externalities},
Proc. EC (2017), 351--368.
\href{https://arxiv.org/abs/1609.06825}{arXiv:1609.06825}.
\bibitem{cast}
M. Castiglioni, A. Celli, A. Marchesi and N. Gatti,
\emph{Online Bayesian Persuasion}, Advances in Neural Information
Processing Systems 33 (2020), 16188--16198.
\href{https://papers.neurips.cc/paper_files/paper/2020/hash/ba5451d3c91a0f982f103cdbe249bc78-Abstract.html}{NeurIPS publication}.
\end{thebibliography}
\end{document}
